\section{The first initialized simulation}\label{sec:sa-histories}
We now use the moving-sheet construction to prove the first result.
Records interspersed among work symbols retain the information erased by
a machine instruction. The processor below first supplies its exact
initialized simulation and the stronger separation of its full rule
cylinders. We then place those cylinders and check the fixed observer
during loading and throughout every subsequent motion.

\subsection{History interspersed among work symbols}\label{subsec:sa-interleaved}
Let $\Gamma$ be the work alphabet, $Q$ the states, and $H$ the halt set.
Adjoin an origin bit to each work symbol, preserving it under writes;
initially only cell zero has bit one.  Write $A=\Gamma\times\{0,1\}$ and
$b_0=(\square,0)$.  For each nonhalt instruction
$\widehat\delta(q,a)=(q',b,d)$, $d\in\{0,+,-\}$, and incoming mode
$h\in\{0,+,-\}$, introduce a distinct record $g_{q,h,a}$.  The records
form an alphabet $G$ disjoint from $A$.  Main controls are $M(q,h)$;
$C_+(q),C_-(q)$ finish a head movement.

\begin{lemma}[Interspersed-history simulation]\label{lem:sa-interleaved}
Starting at $M(q_0,0)$ with left stack $b_0^\infty$ and right stack the
origin-tagged input followed by blanks, this table simulates every ordinary
step in finitely many positive rule applications.  Erasing records at a
main control gives the work tape, head, and state.  The origin bit gives
absolute tape positions.  The full source cylinders are pairwise disjoint,
as are the full image cylinders.  The initialized run reaches a halting
main control exactly when the machine halts.
\end{lemma}
\begin{proof}
At a main control, erasing records makes $R$ the head-and-right tape and
$L$ the reversed left tape.  The first row writes $b$ and saves the lost
information in $g$.  A stay is complete.  For a right move, $b$ has been
put on $L$ and $C_+$ transfers the intervening records until the next
work symbol is exposed.  For a left move, $C_-$ transfers the intervening
records from $L$ and finally its first work symbol; that symbol becomes
the new head immediately before $b$ after erasure.  At finite rule time
there are finitely many records and infinitely many work cells in each
direction, so every scan finishes.  Writes preserve the origin bit.

The complete rule list is
\begin{equation}\label{eq:sa-interleaved-table}
\begin{aligned}
 (M(q,h);\epsilon,a)&\to
 \begin{cases}
 (M(q',0);g,b),&d=0,\\
 (C_+(q');bg,\epsilon),&d=+,\\
 (C_-(q');\epsilon,bg),&d=-,
 \end{cases}\quad g=g_{q,h,a},\\
 (C_+(q);\epsilon,g)&\to(C_+(q);g,\epsilon),&g\in G,\\
 (C_+(q);\epsilon,c)&\to(M(q,+);\epsilon,c),&c\in A,\\
 (C_-(q);g,\epsilon)&\to(C_-(q);\epsilon,g),&g\in G,\\
 (C_-(q);c,\epsilon)&\to(M(q,-);\epsilon,c),&c\in A.
\end{aligned}
\end{equation}
No rule leaves a main control with $q\in H$.

Source separation uses the tested top symbol on $R$ for main and $C_+$
controls, and on $L$ for $C_-$.  Into $C_+$ the left prefix starts with a
work symbol in a main entry and a record in a loop; the record names
separate the main entries.  Into $C_-$ the same distinction is on $R$.
Into $M(q,0)$ the first left record identifies the incoming instruction;
into $M(q,+)$ or $M(q,-)$ the first right work symbol identifies the exit.
Thus every image has a unique inverse on its full cylinder.  The simulation
invariant proves the halt equivalence, including an initial halt.
\end{proof}

We now prove Theorem~\ref{thm:sa-interleaved}.

\begin{proof}
For $K=|A\sqcup G|$, use digits $1,3,\ldots,2K-1$ and base $B=2K+2$.
Enumerate all $m$ controls by $i=1,\ldots,m$, put $\ell=1/[16(m+1)]$,
and use the code
\[
 (\beta_i+\ell E(L),1/4+\ell E(R),1/4),\qquad
 \beta_i=2i\ell+\begin{cases}5/8,&\text{halting},\\1/4,&\text{otherwise}.
 \end{cases}
\]
Nonhalt intervals lie in $(1/4,3/8)$ and halt intervals in $(5/8,3/4)$.
All prefixes in \eqref{eq:sa-interleaved-table} have length at most two;
Lemma~\ref{lem:sa-radix} therefore gives source and target gaps at least
$\ell/B^2$.  Let $N$ be the number of rules.  Apply Lemma~\ref{lem:sa-sheet-routing}, with collar at most
$\min(\ell/(4B^2),1/[16(N+1)],1/16)$, to their full closed rectangles.
Its normal scale compensates the unequal total prefix lengths.
The initial code is rational.  A localized curl translation of radius
$1/32$ loads it along the segment from the fixed label during $[0,1]$.
Repeat the unit rule pulse thereafter.  All pulses have zero collars
inside the cube and are curls, giving a smooth mean-zero velocity $U$
with $U(0)=0$.  Lemma~\ref{lem:sa-realization} supplies the force.

At time $1+k$ induction gives the code after $k$ rules until the first
halt.  A halt is in the observed arc, including at $t=1$ if initial.
For a nonhalting run the loader stays below $1/2$, and every subsequent
horizontal coordinate is constant during vertical motion and between its
two nonhalt endpoints during the middle stage.  Thus the observer is
avoided at every time.  At rule time $k$ the raw stacks are blank beyond
$\max(1,|w|)+2k$ symbols, so this code retains the whole configuration,
not merely the event bit.
\end{proof}

